% \section{Approach}\label{sec:approach}

% Timeline of decisions across a pre-event day and an event day, one swimlane
% per controller. Complements Fig. 1 (dependency map): this figure shows WHEN
% each policy acts and WHICH decision it takes. Lanes are independent clocks;
% arrival and control epochs interleave (no fixed order).
% Requires \usetikzlibrary{arrows.meta, positioning, backgrounds}.
\begin{figure*}[t]
\centering
\begin{tikzpicture}[
  font=\scriptsize,
  lanelab/.style={align=center, anchor=east, font=\scriptsize},
  evt/.style={font=\tiny\itshape, text=black!65, align=center},
  arr/.style={-{Latex[length=1.6mm]}, thick},
  x=1cm, y=1cm
]
% day d-1: x in [0,6.4]; overnight: [6.4,8.0]; day d: [8.0,14.4]
% lanes: DR y=2.5, ch y=1.25, neg y=0
% ---------- event window band (day d, behind everything) ----------
\begin{scope}[on background layer]
\fill[red!8] (11.4,-0.75) rectangle (13.3,3.05);
\draw[red!40, thin] (11.4,-0.75) rectangle (13.3,3.05);
\end{scope}
\node[evt, text=red!60!black] at (12.35,3.25) {event window $\evwin{d}$ (4--5\,h)};
% ---------- day headers ----------
\node[font=\scriptsize\bfseries] at (3.2,3.6) {day $d{-}1$ (pre-event)};
\node[font=\scriptsize\bfseries, text=black!55] at (7.2,3.6) {overnight};
\node[font=\scriptsize\bfseries] at (9.9,3.6) {day $d$ (event called)};
\draw[black!30, thin] (0,3.4) -- (6.4,3.4);
\draw[black!30, thin, dashed] (6.4,3.4) -- (8.0,3.4);
\draw[black!30, thin] (8.0,3.4) -- (14.4,3.4);
% ---------- lane baselines ----------
\draw[blue!35, thick] (0,2.5) -- (6.4,2.5);
\draw[blue!35, thick, densely dotted] (6.4,2.5) -- (8.0,2.5);
\draw[blue!35, thick] (8.0,2.5) -- (14.4,2.5);
\draw[teal!45, thick] (0,1.25) -- (6.4,1.25);
\draw[teal!45, thick, densely dotted] (6.4,1.25) -- (8.0,1.25);
\draw[teal!45, thick] (8.0,1.25) -- (14.4,1.25);
\draw[orange!55, thick] (0,0) -- (6.4,0);
\draw[orange!55, thick, densely dotted] (6.4,0) -- (8.0,0);
\draw[orange!55, thick] (8.0,0) -- (14.4,0);
% ---------- lane labels ----------
\node[lanelab, text=blue!50!black]  at (-0.25,2.5)
  {$\pidr$ \\ {\tiny once per day}};
\node[lanelab, text=teal!50!black]  at (-0.25,1.25)
  {$\pich$ \\ {\tiny every $\dt$}};
\node[lanelab, text=orange!60!black] at (-0.25,0)
  {$\pineg$ \\ {\tiny per arrival,} \\ {\tiny departure}};
% ---------- pi^DR: start-of-day commitments ----------
\node[diamond, fill=blue!25, draw=blue!60, inner sep=1.6pt] at (0,2.5) {};
\node[evt, anchor=south west, text width=3.2cm] at (-0.05,2.64)
  {commit $\FSL_d$ for day $d$ (Alg.~\ref{alg:dr})};
\draw[arr, blue!55] (0.15,2.42) -- (0.85,1.40)
  node[evt, pos=0.75, right=1pt, text width=1.6cm] {targets steer current SoC banking};
\node[diamond, fill=blue!25, draw=blue!60, inner sep=1.6pt] at (8.0,2.5) {};
\node[evt, anchor=south west, text width=1.9cm] at (7.95,2.64)
  {commit $\FSL_{d+1}$};
\draw[arr, blue!55] (8.6,2.42) -- (8.6,1.40)
  node[evt, pos=0.45, right=1pt, text width=1.5cm] {$\FSL_d$ in effect for day $d$};
% ---------- pi^ch: control comb ----------
\foreach \xx in {0.0,0.32,...,6.4}  \draw[teal!60] (\xx,1.18) -- (\xx,1.32);
\foreach \xx in {8.0,8.32,...,14.4} \draw[teal!60] (\xx,1.18) -- (\xx,1.32);
\node[evt, anchor=north, text width=3.6cm] at (3.0,1.11)
  {Charge excess SoC when ($\Ptot{t}<\Pmaxhat$), pulled to targets (Alg.~\ref{alg:ch})};
\node[evt, anchor=north, text width=2.2cm] at (12.35,1.11)
  {discharge buffer; hold $\Ptot{t}\le\FSL_d$};
% ---------- pi^neg: arrivals and departures ----------
\foreach \xx in {1.1,1.7,2.4}
  \draw[arr, orange!75!black] (\xx,0.32) -- (\xx,0.04);
\foreach \xx in {9.1,9.7,10.4}
  \draw[arr, orange!75!black] (\xx,0.32) -- (\xx,0.04);
\foreach \xx in {4.9,5.5,6.1}
  \draw[arr, orange!75!black] (\xx,-0.04) -- (\xx,-0.32);
\foreach \xx in {13.5,13.9,14.3}
  \draw[arr, orange!75!black] (\xx,-0.04) -- (\xx,-0.32);
\node[evt, anchor=north, text width=3.2cm] at (1.75,-0.42)
  {arrivals: flexibility menu, banking consent at $\pbank{v}{k}$};
\node[evt, anchor=north, text width=3.1cm] at (5.5,-0.42)
  {departures};
\node[evt, anchor=north, text width=3.0cm] at (13.4,-0.85)
  {departures $\Ubuf$, $\Udc$};
% ---------- overnight carryover ----------
\draw[arr, black!60, densely dashed] (6.5,0.62) -- (7.9,0.62);
\node[evt, text width=1.7cm] at (7.2,1.12)
  {Excess SoC carryover};
\end{tikzpicture}
\caption{When each decision is taken, and by which controller, across a
pre-event day and an event day. The commitment policy $\pidr$ acts once per
day, at the start of the day: at the start of day $d{-}1$ it commits
$\FSL_d$ for the following day, before any day-$(d{-}1)$ or day-$d$ vehicle
has arrived, and sets the terminal targets that steer the current day's
banking. The charging policy $\pich$ acts every control step: on the
pre-event day it banks excess SoC into headroom, and inside the called
window it discharges the surviving buffer to hold the committed level. The
negotiation policy $\pineg$ acts at each arrival (flexibility menu and
banking consent) and each departure (settlement on realized quantities; on
event days, the replay shares $\Ubuf$ and $\Udc$). The three rows are
independent clocks whose epochs interleave; the commitment is fixed a full
day before the day it governs (Section~\ref{sec:coupled}).}
\label{fig:timeline}
\end{figure*}

% Section~\ref{sec:formulation} posed three sub-problems and Section~\ref{sec:coupled} showed they compose into one semi-Markov program. This section describes how each policy is computed and how the three exchange information at run time. 





\subsection{Day-Ahead Commitment ($\pidr$)}\label{sec:approach-dr}



The FSL commitment policy chooses $\FSL_d$ at the start of day $d{-}n$ by solving a day-ahead problem while evaluating its risk. It solves a MILP that minimizes expected cost over sampled days drawn from the forecast $\psihat$, providing a possible event-window peak. 
% This is the Sample Average Approximation (SAA) over the sampled future states. where the true violation probability is replaced by its empirical frequency over the sampled days. The forecast $\psihat$ already includes the SoC banking and flexibility the building anticipates, enabling the nomination to reflect the resources it expects to have. The negotiation of Section~\ref{sec:approach-neg} then realizes that flexibility on the day and settles it as revenue and user savings, but it never revises the committed level.

\begin{algorithm}[t]
\caption{$\pidr$: day-ahead commitment by scenario MILP}
\label{alg:dr}
\DontPrintSemicolon
\SetCommentSty{algocmt}
\KwIn{forecast $\psihat$ (including anticipated banking and flexibility);
event window $\evwin{d}$; delivery risk level $\varepsilon$; commitment
depth $\lambda \in [0,1]$; sample count $N$; baseline peak $\Phist$;
per-vehicle usable maxima $\{E_{\max}(v)\}$}
\KwOut{committed level $\FSL_d$}
set $\etarget{v}{k} \leftarrow E_{\max}(v)$ for every user $v$ and session $k$
\tcp*{uniform \emph{aspiration}; clipped by the headroom gate in $\pich$}
Draw $N$ sampled days $\omega_1,\dots,\omega_N \sim \psihat$ (arrivals,
required SoC, off-site use)\;
solve the day-ahead scenario MILP once over $\omega_1,\dots,\omega_N$,
minimizing expected cost \tcp*{risk-neutral nomination}
\lFor{each $\omega_i$}{$m_i \leftarrow \max_{t \in \evwin{d}} \Ptot{t}$
\tcp*[f]{window peak realized by that solve}}
$\min\Omeps \leftarrow$ the $(1-\varepsilon)$ empirical quantile of
$\{m_1,\dots,m_N\}$ \tcp*{deepest cap held on $\ge 1-\varepsilon$ of days}
$\FSL_d \leftarrow \Phist - \lambda\,\bigl(\Phist - \min \Omeps\bigr)$
\tcp*{commitment-depth interpolation}
\Return $\FSL_d,\ \{\etarget{v}{k}\}$\;
\end{algorithm}

% A cap $F$ is held on sampled day $i$ exactly when that day's window peak $m_i = \max_{t\in\evwin{d}} \Ptot{t}$ stays at or below it: \begin{equation}\label{eq:xflag} \text{day $i$ holds cap $F$} \;\Longleftrightarrow\; m_i \le F \;\Longleftrightarrow\; \Ptot{t} \le F \ \text{for every } t \in \evwin{d}. \end{equation} The commitment is thus evaluated \emph{jointly} over the window: a day counts as a failure if it breaches the cap at \emph{any} within-window step, not as a per-step average. A level that passes most steps can still fail the joint test, and only the joint test matches the delivery guarantee the building sells. Taking the $(1-\varepsilon)$ quantile of the peaks $\{m_i\}$ therefore returns $\min\Omeps$, the deepest cap held on at least a $1-\varepsilon$ fraction of sampled days. Because the peaks $m_i$ are read from the risk-neutral solve, which was not itself trying to reach deeper caps, $\min\Omeps$ is a \emph{conservative} estimate of the deepest deliverable level; a per-candidate delivery pass would only lower it, so the reported commitments understate rather than overstate what the fleet can hold. The interpolation returns the null commitment $\FSL_d = \Phist$ at $\lambda = 0$ and this deepest $\varepsilon$-feasible level at $\lambda = 1$, and Proposition~\ref{prop:newsvendor} identifies the risk-neutral anchor within this family. 

The terminal targets $\etarget{v}{k}$ set how much buffer the commitment layer \emph{asks} each user to leave behind; they are a soft pull, not a hard requirement. The current implementation uses a uniform aspiration, $\etarget{v}{k} = E_{\max}(v)$, nominally asking every vehicle to bank to its usable maximum. This does not mean vehicles are charged to full: the target only biases dispatch, and the banking it induces is admitted by $\pich$ only into headroom below the estimated and running peaks (Section~\ref{sec:approach-ch}) and always within the committed FSL and rate limits. A vehicle therefore reaches $E_{\max}$ only if it can do so without violating the estimated peak $\Pmaxhat$. 

% The uniform aspiration is deliberately simple and, being maximal, un-prioritized: it does not shape the buffer to the commitment's actual shortfall or choose which vehicles should carry it. A refinement that distributes only the buffer the commitment requires (from $\min\Omeps$), across users in proportion to their forecast availability, would reduce unnecessary banking and its degradation cost without changing the interface to $\pich$; we treat it as future work and report results under the uniform aspiration.

\subsection{Real-Time Charging ($\pich$)}\label{sec:approach-ch}

The charging policy is a Monte Carlo model-predictive controller (MC-MPC): it acts every control step under a receding horizon, and handles the uncertainty in future arrivals by sampling. At step $t$ it draws $K$ samples from $\psihat$, solves a mixed-integer program over the remaining horizon for each with a common first action and the committed $\FSL_d$. The first action at the current time $t$ is applied, and rolls forward. The objective is Eq.~\eqref{eq:chobj}, combining energy, demand, the FSL penalty, degradation, and the end-of-day SoC banking reward $\Kbank{v}\bshort{v}$ that pulls departures toward $\etarget{v}{k}$.


\begin{algorithm}[tbph]
\caption{P-V2B Charging Optimization}
\label{alg:MP-MPC}
\KwIn{Initial state $S_0$, number of samples $N_f$, billing period $T$, estimated monthly peak power $\hat{P}^{max}$ 

}
\KwOut{Charging decisions $A_t$}

\For{$t = 0$ \KwTo $\mathcal{T}$}{

    Check EV departures and free assigned chargers\;
    Assign chargers to new arrivals using FCFS\; 
    Observe current state $\State_t$\;
    Future trajectory set $\mathcal{F} \leftarrow \emptyset$\; 
    \For{$i = 1$ \KwTo $N_f$}{
       Sample trajectory $f$ from the generative model given $\State_t$, simulate from $t{+}1$ to $t_{eod}$, and add $f$ to $\mathcal{F}$\;
    }
    Get $A_t$ by solving optimization in Eq.~(\ref{eq:ns_objective}) 
    
    Refine action to $\tilde{A}_t=f_{\text{refine}}(A_t, \hat{P}^{\max})$ by Eq.~(\ref{eq:force_charge})  
    % Subject to: $A_t$ is shared across all $f$, and feasibility constraints (Section~\ref{sec:constraints})\;
    
    Apply $\tilde{A}_t$ and update system state to $\State_{t+1}$\;
}
% \vspace{-1mm}
\end{algorithm}

Two components govern SoC banking. The target SoC $\etarget{v}{k}$, through $\Kbank{v}\bshort{v}$, states the size of the buffer and the  estimated peak $\Pmaxhat$ and the running cycle peak $\pmaxpast$ estimated cap $\Pmaxhat$ states when accumulating it is favorable. Hence, banking never creates a new demand-charge peak. When no event is anticipated, $\muDR = 0$ and $\Kbank{v} = 0$, and the composition reduces exactly to the opportunistic headroom rule of P-V2B, which is the null-commitment special case. 


% The dual $\muDR$ on the FSL-delivery requirement is the shadow price of that requirement: the marginal improvement in the building's objective per additional unit of guaranteed deliverable buffer, and hence the value the banking price passes to the user. A banking offer is quoted at arrival, possibly days before an event, so $\muDR$ at quote time cannot come from that day's operation; it is the \emph{expected} shadow price of delivering on the anticipated event, evaluated under $\psihat$. It is zero when no event is anticipated and positive on pre-event days in proportion to the buffer the building forecasts it will need. Being a forecast quantity, $\muDR$ is the single price-forecast entering any user-facing offer, consistent with the incentive analysis of Section~\ref{sec:guarantees}.


\subsection{Per-Arrival Negotiation and Settlement ($\pineg$)}\label{sec:approach-neg}

The negotiation policy acts at each arrival and each departure. On arrival it prices the flexibility menu and the banking option under the committed-FSL objective and presents them alongside the reject-all outside option; on departure it settles every user-facing term on realized quantities through a single replay of the day.  The banking price makes the over-charge free exactly when the building values the buffer, through the $\muDR\rhohat_v$ term, and arbitrage-neutral otherwise, through the $\max$; the derivation and its incentive consequence are in Section~\ref{sec:formulation} and Section~\ref{sec:guarantees}. The departure settlement re-solves the realized day with rates as the only free variables, so the difference between the with-buffer and without-buffer solves is the true marginal value of $v$'s buffer; partitioning that one difference into $\Udc$ and $\Ubuf$ attributes the demand-charge and the energy-plus-penalty shares without paying any step twice, and the proportional cap of Lemma~\ref{lem:replaybf} keeps the total within the realized saving.

\subsection{Coupling and Complexity}\label{sec:approach-coupling}

The four coupling quantities are the entire interface between policies.
Within a day the dependence is one-way: $\FSL_d$ and $\etarget{v}{k}$
constrain and steer $\pich$, whose realized dispatch and dual $\muDR$
set what $\pineg$ can offer; nothing flows back within the day. Across days only the surviving buffer $\buff{d{+}1}$ and the running peak enter the next commitment. The terminal targets $\{\etarget{v}{k}\}$ are precisely the per-user, per-session scalars of Theorem~\ref{thm:decomp}'s policy class: carrying the commitment down into charging through this small set of numbers is what makes the per-day decomposition exact.
The dominant cost is the model-predictive charging solve: one
mixed-integer program per control step per sampled continuation, and
$N$ delivery simulations per candidate level in the day-ahead pass. Both
are parallel across sampled days, and the day-ahead pass is amortized over
the whole operating day; the settlement replay adds two rates-only
solves per departing user, reusing the same delivery model.